% www.electronics-tutorials.ws cute, AND and OR

\subsection{\python\ come linguaggio interpretato}

Bytecode come primo passo dell'esecuzione di un programma e cartella pycache.
Il file intermedio costituisce una serie di istruzioni per la Python Virtual Machine (PVM).
Il bytecode \`e ottimizzato per l'esecuzione, ed il risultato \`e pi\`u veloce che
non interpretare il codice sorgente direttamente.

Vogliamo mettere un accenno a \url{https://docs.python.org/3/library/dis.html}?

Confrontato con un tipico linguaggio compilato, come ad esempio il C, un linguaggio
interpretato come \python\ fornisce indubbiamente, dal punto di vista del programmatore,
una serie di vantaggi da non sottovalutare. In C il tipo delle variabili deve essere
esplicitamente dichiarato, e durante l'esecuzione di un programma si pu\`o cambiare
il valore di una variabile ma non il tipo; in \python\ non \`e necessario dichiarare
le variabili, e non ci sono vincoli di sorta sui tipi. In Python \`e estremamente
pi\`u semplice create container (e.g., liste) eterogenei---cio\`e contenenti oggetti
di tipo diverso. In C la gestione della memoria \`e manuale, mentre in \python\
\`e automatica attraverso un meccanismo di garbage collection. In generale la sintassi
di \python\ \`e considerata dalla maggior parte dei programmatori come chiara, espressiva
e concisa. Ma se tutto questo \`e vero, perch\'e esistono ancora i linguaggi compilati?

Beh, tutte queste caratteristiche desiderabili di \python\ come linguaggio interpretato,
come del resto tutto nella vita, non vengono a costo zero. Il fatto che non dobbiamo
dichiarare le variabili, che possiamo cambiare il tipo a runtime, e che non dobbiamo
preoccuparci di liberare la memoria quando non ci serve pi\`u fa s\`i che anche
gli oggetti pi\`u semplici (e. g., un numero intero) in \python\ debbano tener traccia
autonomamente di un certo numero di cose, e. g., di che tipo sono e quante referenze
hanno; la conseguenza \`e che il footprint in memoria di un programma \`e tipicamente
molto pi\`u grande in \python\ che in C. (Vedremo questo pi\`u in dettaglio nella
sezione~\ref{sec:rappresentazione_interi}.) Il fatto che una lista in \python\ sia
intrinsecamente un contenitore eterogeneo rende l'iterazione lenta---ed a volte
\emph{estremamente} lenta. (Anche se, come vedremo nella sezione~\ref{sec:numpy}
ci sono espedienti per ovviare a questo.)

In un mondo in cui i computer sono veloci e la memoria a disposizione \`e copiosa,
in molti casi queste non sono limitazioni rilevanti, ed \`e spesso vantaggioso rendere
le cose pi\`u facili per la persona che sviluppa il codice piuttosto che per il
computer che lo esegue. Ma \`e bene avere chiaro fin da subito che vi sono innumerevoli
situazioni in cui questo trade-off si inverte, ed allora \`e necessario fare compromessi.


\subsection{Controllo del flusso}

"Ripetete il tutto per altre tre volte."; "lasciate cuocere fino a che le patate
non saranno dorate"; "se il sugo \`e troppo secco aggiungere un po' di acqua di
cottura". Se avete letto anche distrattamente un libro di ricette in vita vostra
vi sar\`a capitato di imbattervi in frasi di questo tipo. Per noi le parole chiave
sono \emph{tre volte}, \emph{fino a che}, e \emph{se}.

Una delle cose che capita pi\`u di frequente quando programmiamo \`e quella di
\emph{ripetere}, o \emph{iterare} un'operazione. Sommare i numeri interi da $0$ a
$100$, per esempio, significa in pratica partire da $0$ e poi sommare ripetutamente
$1$, $2$\ldots e cos\`i via. \python, cos\`i come del resto la maggior parte dei
linguaggi di programmazione offre due costrutti: \code{for} e \code{while}.




\subsection{Imparare ad amare il \emph{backslash}}

%https://lerner.co.il/2018/07/24/avoiding-windows-backslash-problems-with-pythons-raw-strings/
Invariabilmente il neofita che inizia ad usare \python\ sotto \windows\ per leggere un
file (diciamo, per fissare le idee, che il percorso a quest'ultimo \`e
\texttt{C:$\backslash$abc$\backslash$def$\backslash$ghi.txt}) si trova a dover affrontare
la situazione, apparentemente inspiegabile, per cui \python\ si ostina a dire che
il file non esiste. Eppure abbiamo controllato innumerevoli volte, ed il file \`e l\`i.
Cerchiamo di capire che cosa sta succedendo, iniziando dal frammento~\ref{snip:dos_path_escape}.

\snip{dos_path_escape}{
Caption.
}

Qualcosa chiaramente non va: abbiamo una \cchar{a} ed un backslash mancanti all'appello
dopo i due punti, nella stringa di testo in uscita. La questione \`e che il carattere
\cchar{\textbackslash} (ovvero il backslash) ha un significato particolare in molti
linguaggi di programmazione, incluso \python: conferisce un significato differente
dal solito al carattere (o alla serie di caratteri) che lo segue, in quello che
tecnicamente si chiama una sequenza di escape, e che tipicamente si usa per rappresentare
i caratteri di controllo della tabella ASCII (codici $0$--$31$). L'esempio forse
pi\`u noto \`e quello dell'andata a capo, o line feed, che corrisponde al codice di
controllo $10$ della tabella ASCII, e che nella quasi totalit\`a dei linguaggi di
programmazione si rappresenta come \cchar{\textbackslash{}n}, ma \python, con una
scelta mutuata dal C, riconosce, tra le altre, le sequenze di escape nella
tabella~\ref{tab:sequenze_escape}.

\begin{table}[!htb]
  \tablehstack{
    \centering\begin{tabular}{lll}
        \hline
        Sequenza di escape & Significato & Codice ASCII\\
        \hline
        \hline
        \textbackslash{}\textbackslash & Backslash (\cchar{\textbackslash}) & \\
        \textbackslash{}' & Single quote (\cchar{'}) & \\
        \textbackslash{}" & Double quote (\cchar{"}) & \\
        \textbackslash{}a & ASCII Bell (BEL) & \\
        \textbackslash{}b & ASCII Backspace (BS) & \\
        \textbackslash{}f & ASCII Formfeed (FF) & \\
        \textbackslash{}n & ASCII Linefeed (LF) & $10$\\
        \textbackslash{}r & ASCII Carriage Return (CR) & \\
        \textbackslash{}t & ASCII Horizontal Tab (TAB) & \\
        \textbackslash{}v & ASCII Vertical Tab (VT) & \\
        \hline
      \end{tabular}
  }{
    \caption{Tabella delle sequenze di escape riconosciute da \python, come illustrato nella
      \href{https://docs.python.org/3/reference/lexical_analysis.html\#escape-sequences}{documentazione}.
      Sebbene le sequenze di escape non siano esattamente la cosa che incontrerete
      pi\`u spesso nella vita (\cchar{\textbackslash{}n} a parte per andare a capo),
      il TAB (\cchar{\textbackslash{}t}) \`e occasionalmente utile per allineare
      colonne di numeri all'interno di file di testo.}
    \label{tab:sequenze_escape}
  }
\end{table}

Cominciamo lentamente a capire cosa sta succedendo. Nel frammento~\ref{snip:dos_path_escape}
i due caratteri \cchar{\textbackslash{}a} sono interpretati (correttamente) come la
sequenza di escape corrispondente al codice di controllo $7$ nella tabella ASCII,
vale a dire BEL, ovverosia la campanella del terminale---se eseguite il frammento
interattivamente nel prompt di \python\ e non avete disabilitato la campanella del
terminale, dovreste sentire un suono!

Tornando al problema del percorso al nostro file di dati, l'altro aspetto rilevante
della questione \`e la scelta degli sviluppatori della versione 2.0 del DOS di scegliere
il backslash come separatore delle directory nel filesystem~\cite{dos_backslash},
scelta che si \`e propagata a \windows\ ed \`e arrivata fino ai nostri giorni\footnote{
I sistemi Unix e Unix-like, incluso GNU-Linux e Mac-OS, utilizzano uno \cchar{/}
(slash) come carattere per la separazione delle directory nel filesystem, e questo
problema, molto semplicemente, non si pone.}.
Solo che in una stringa di \python (e, per quel che conta, lo stesso vale per la
maggior parte dei linguaggi di programmazione e dei terminali) il backslash ha un
significato speciale, e non \`e trattato come tutti gli altri caratteri. Allora la
domanda \`e: come faccio a scrivere un backslash in una stringa? Una prima risposta
corretta ce la fornisce la prima riga della tabella~\ref{tab:sequenze_escape}: possiamo
usare un backslash per fare l'escape di se stesso, il che vuol dire che la sequenza
di escape \cchar{\textbackslash{}\textbackslash} si traduce in un \cchar{\textbackslash{}}
semplice. Questo funziona come uno si aspetta nella maggior parte delle situazioni.
In alternativa, \python\ fornisce la possibilit\`a di aggiungere una lettera \cchar{r}
o \cchar{R} prima della virgoletta iniziale della stringa: cos\`i facendo si crea
quella che si chiama una stringa grezza (raw string), in cui il backslash \`e trattato
alla pari di un qualsiasi altro carattere. Il frammento~\ref{snip:dos_path_escape_fix}
illustra le due possibilit\`a.

\snip{dos_path_escape_fix}{
Caption.
}

C\`e un ultimo, piccolo colpo di scena nella storia, ovvero il fatto che, come vederemo
nella sezione~\ref{sec:latex}, il backslash \`e utilizzato in \LaTeX\ per identificare
i comandi all'interno di un documento. (Questo crea tutta una serie di sottili difficolt\`a
quando si cerca di scrivere un backslash in un documento \LaTeX, e non vi nascondo
che la battitura di questa sezione \`e stata meno triviale di quanto mi sarei aspettato).





\subsection{Digressione: stringhe, numeri e modalit\`a di IO}

Arrivati a questo punto, se avete letto con attenzione, non potrete fare a meno di
chiedervi: se io scrivo su disco un singolo byte, ad esempio \byte{10101001},
che cosa ho fatto realmente? Dato che
\begin{align*}
  0b10101001 = 0xA9 = 169
\end{align*}
abbiamo almento tre risposte potenzialmente corrette a disposizione! Il $169$ \`e
al di fuori della tabella ASCII ristretta, ma se lo interpretiamo come code point
di una tipica estensione ad $8$ bit della codifica ASCII, abbiamo scritto un bel
simbolo di copyright. Se assumiamo che il nostro byte di storage rappresenti un
intero unsigned, allora abbiamo scritto $169$. Se invece assumiamo che sia un intero
con segno, siccome il bit pi\`u significativo \`e $1$, si tratta di un numero negativo,
e per la precisione $-87$. Complicato, eh?

La risposta corretta \`e: dipende. Quando rileggiamo dal disco il nostro byte di
informazione possiamo intepretarlo come vogliamo. Ed il corollario \`e che sar\`a
bene assicurarci prima dell'intento con cui il file \`e stato scritto per evitare
di prendere cantonate. Se il file in questione \`e un file di testo, allora dovremo
leggerlo come tale, e con l'encoding corretto (quello con cui \`e stato scritto).
Se invece il file \`e un file binario, allora dovremo leggerlo con la modalit\`a
opportuna, e dobbiamo sapere esattamente qual \`e il layout dei dati all'interno
per poterli spacchettare nel modo opportuno.

\snip{byte}{
Caption.
}

Abbiamo imparato una lezione fondamentale, ovvero la distinzione tra il contenuto
di una porzione di memoria (sia essa volatile o su disco) ed il significato che
noi diamo a quella particolare sequenza di $0$ ed $1$ all'interno del nostro programma.
Ed abbiamo capito che lo stesso bit pattern in memoria pu\`o essere interpretato
in pi\`u modi, che che dobbiamo fare attenzione ad essere consistenti nelle fasi di
lettura e di scrittura.

Prima di chiudere vale la pena di parlare un attimo della scrittura su disco. Senza
entrare in discorsi troppo complicati, il caso d'uso tipico per il corso di laboratorio
\`e quello in cui scriviamo, manualmente o programmaticamente, un numero (limitato)
di valori in un file di testo, e poi rileggiamo questi valori con \numpy\ per
analizzarli, come spiegato sommariamente nella sezione~\ref{sec:numpy_loadtxt}.
Questo va benissmo, e funziona come uno si aspetterebbe, ma non dobbiamo dimenticare
che, mentre un numero in virgola mobile (lo vedremo nella sezione~\ref{sec:virgola_mobile})
costa su disco $32$ o $64$~bit, un numero scritto in un file di testo costa $8$
bit per ciascun carattere---incluse le cifre numeriche, il separatore decimale se
presenti, e la potenza di $10$ se scriviamo in notazione ingegneristica. Tanto per
fissare le idee, scrivere \texttt{9.1093837139e-31} (la massa dell'elettrone in kg)
in un file di testo costa esattamente $128$~bit, mentre con la met\`a dello spazio
potremmo confortevolmente raggiungere lo stesso risultato se scrivessimo un file
binario. \`E probabile che non abbiate l'occasione di apprezzare l'importanza della
cosa quest'anno, ma tenetelo a mento, perch\'e prima o poi succeder\`a.




Siamo finalmente arrivati alla parte interessante. Lo standard IEEE 754~\cite{ieee_754},
che è ormai adottato pressoché universalmente, fornisce una prescrizione
per la rappresentazione dei numeri in virgola mobile in logica binaria, e
due specifiche di formato: a 32~bit (in precisione singola) e 64~bit (in
precisione doppia), come mostrato in figura~\ref{fig:rappresentazione_float}.
Più precisamente, ogni volta che scriviamo un numero in virgola mobile, dobbiamo
pensare che, all'interno del calcolatore, esso è rappresentato come
\begin{align}\label{eq:ieee_754}
  x = (-1)^s \times m \times 2^{e - b} \quad \text{dove} \quad 1 \leq m < 2,
\end{align}
in cui, spostandosi da sinistra a destra nella rappresentazione---ovverosia dal
bit più significativo a quello meno significativo:
\begin{itemize}
  \item $s$ rappresenta il \emph{segno} del numero, e corrisponde al bit
    più significativo ($0$ se il numero è positivo, $1$ se è negativo);
  \item $e$ è l'esponente, che ha a disposizione 8~bit in precisione singola
    e 11~bit in precisione doppia; $b$, che prende il nome di \foreign{bias},
    è una costante additiva (positiva) che permette di rappresentare numeri
    $< 1$ senza bisogno di aggiungere un segno all'esponente, e vale $127$ in
    precisione singola e $1023$ in precisione doppia;
  \item $m$ prende il nome di \emph{mantissa} e rappresenta (utilizzando 23~bit in
    precisione singola e 52~bit in precisione doppia) la sequenza di cifre dopo
    la virgola, \emph{assumendo che la parte intera sia $1$} (cioè la mantissa
    è convenzionalmente compresa tra $1$ e $2$ ed il bit più significativo,
    che è garantito essere $1$, può essere omesso).
\end{itemize}

\begin{figure}[!htbp]
  \center\input{figures/rappresentazione_float}
  \vspace*{-10pt}
  \caption{\foreign{Layout} dei due formati per la rappresentazione binaria dei
  numeri in virgola mobile definiti dallo standard IEEE~754: in precisione
  doppia (11~bit per l'esponente e 53 per la mantissa) ed in precisione singola
  (8~bit per l'esponente e 23 per la mantissa). Il \foreign{bias} per l'esponente
  è $1023$ nel primo caso e $127$ nel secondo.}
  \label{fig:rappresentazione_float}
\end{figure}


Cominciamo allora da un numero semplice (i.e., con uno sviluppo binario finito,
come quello che abbiamo considerato all'inizio di questa sezione):
\begin{align*}
  x = 2.25 = 0b10.01 = \frac{9}{4} = 1.125 \times 2.
\end{align*}
I passaggi sono tutti rilevanti (e.g., $\nicefrac{9}{4}$ costituisce la
rappresentazione di $x$ come il rapporto tra i due interi più piccoli possibile,
come potete verificare direttamente utilizzando \pyfunc{as_integer_ratio} in \python)
ma l'ultimo è quello più importante, perché rappresenta $x$ esattamente
nella forma~\eqref{eq:ieee_754}, ovverosia come il prodotto di una mantissa
$1 \leq m < 2$ per una potenza di $2$.
Nel formato in doppia precisione avremo dunque
\begin{align*}
  s & = 0\\
  e & - b = 1 \implies e = 1024 =  0b10000000000\\
  m & = 1.125 = 0b(1).0010000000000000000000000000000000000000000000000000
\end{align*}
(notate che abbiamo indicato tutti gli 11~bit per l'esponente ed i 52~per la
mantissa, con il bit più significativo di normalizzazione tra parentesi) e,
in definitiva
\begin{align*}
  x \rightarrow 0b0\overbrace{0010000000}^{e-b}
  \overbrace{0010000000000000000000000000000000000000000000000000}^{m~\text{(bit più significativo omesso)}}
  = 0x4002000000000000 .
\end{align*}
Wow! Abbiamo scritto il nostro primo numero in virgola mobile nella forma in
cui lo troveremmo nel disco rigido del nostro calcolatore! Adesso possiamo
chiudere il cerchio e provare a fare la stessa cosa con il famigerato
\begin{align*}
  x = 0.1 = 0.1 \times 16 \times 2^{-4} = 1.6 \times 2^{-4}.
\end{align*}
Procedendo esattamente come prima si ha
\begin{align*}
  s & = 0 \\
  e & - b = - 4 \implies e = 1019 = 0b01111111011\\
  m & = 1.6 \approx 0b(1).1001100110011001100110011001100110011001100110011010
\end{align*}
(notate che questa volta la mantissa non può rappresentare esattamente il numero di
partenza con un numero finito di bit) e dunque
\begin{align*}
  x \rightarrow 0b0\overbrace{01111111011}^{e-b}
  \overbrace{1001100110011001100110011001100110011001100110011010}^{m~\text{(bit più significativo omesso)}}
  = 0x3FB999999999999A.
\end{align*}
La cosa interessante, a questo punto, è che se moltiplichiamo la mantissa per
$2^{52}$ (il che equivale a considerare la mantissa stessa, incluso il bit più
significativo \emph{nascosto}, come un numero intero espresso in base $2$), si ha
\begin{align*}
  x = \frac{m \times 2^{52}}{2^{52 - (e  - b)}} =
  \frac{7205759403792794}{2^{56}} = \frac{3602879701896397}{2^{55}}
\end{align*}
(l'ultimo passaggio è possibile grazie al fatto che il numeratore è divisibile
per $2$, in quanto la cifra meno significativa del suo sviluppo binario è zero).
Ricorda qualcosa? Il frammento~\ref{snip:integer_ratio}?


\section{Buone e cattive proprietà dell'aritmetica in virgola mobile}

Adesso che abbiamo gli strumenti di base per una discussione più informata sulle
proprietà dell'aritmetica in virgola mobile su un calcolatore digitale,
facciamo un passo indietro e riesaminiamo criticamente alcune delle cose che
abbiamo detto all'inizio di questa sezione.

Il modello di rappresentazione dei numeri codificato nello standard IEEE~754
ha una serie di buone proprietà che lo rendono estremamente versatile e
riducono, per quanto possibile, il potenziale impatto della precisione finita
sui risultati finali: permette di rappresentare numeri su un intervallo dinamico
di molti ordini di grandezza e fornisce la stessa accuratezza \emph{relativa}
su tutto l'intervallo. In particolare: il massimo numero rappresentabile è dettato
da numero di bit $n_e$ dell'esponente
\begin{align*}
  x_\text{max} \approx 2^{2^{n_e - 1}} =
  \begin{cases}
    2^{128} \approx 10^{38} \quad & \text{in precisione singola}\\
    2^{1024} \approx 10^{308} \quad & \text{in precisione doppia}.
  \end{cases}
\end{align*}
Se si considera che il numero stimato di protoni nell'universo è dell'ordine di
$10^{80}$ e la lunghezza di Planck è $1.6 \times 10^{-35}$~m, capiamo
immediatamente che non c'è di che preoccuparsi.

Di converso, la \emph{precisione} è dettata dal numero $n_m$ di bit a disposizione
per la mantissa, e si può esprimere in termini di cifre decimali equivalenti
come
\begin{align*}
  \frac{x}{\varepsilon_x} \approx \log_{10} \left( 2^{n_s + 1}\right) =
  \begin{cases}
    7 \quad & \text{in precisione singola}\\
    16 \quad & \text{in precisione doppia}.
  \end{cases}
\end{align*}
Anche qui, almeno apparentemente, non abbiamo niente da temere---il momento magnetico
dell'elettrone è misurato in circa una parte su $10^{12}$, e questo
rappresenta più o meno lo stato dell'arte.

\snip{floating_point2}{%
Semplice illustrazione di come la proprietà commutativa dell'addizione non
valga quando applicata all'aritmetica in virgola mobile su un calcolatore
digitale.}

Eppure non possiamo chiudere gli occhi e riposare tranquilli, perché l'aritmetica in
virgola mobile è prodiga di sorprese, come mostrato nel frammento di
codice~\ref{snip:floating_point2}. La referenza~\cite{bush_floating_point}
contiene una discussione più approfondita della questione ed una serie di consigli
pratici per evitare i problemi più comuni (e.g., sottrazione di numeri molto grandi
vicini tra loro, operazioni con numeri che differiscono di molti ordini di grandezza,
confronto tra numeri in virgola mobile nel controllo del flusso logico del programma),
ma adesso è veramente giunto il momento di chiudere.



\subsection{Leggere file di testo con tabelle di numeri}
\label{sec:numpy_loadtxt}

